\documentclass[11pt]{article}
\usepackage[T1]{fontenc}
\usepackage{lmodern}
\usepackage[letterpaper,margin=1in]{geometry}
\usepackage{amsmath,amssymb,amsthm,mathtools}
\usepackage{microtype}
\usepackage[hidelinks]{hyperref}
\usepackage{enumitem}
\hypersetup{pdftitle={Uniform Logarithmic Gamma Coercivity at Fixed Exponential Complexity},pdfauthor={Jongmin Choi},pdfsubject={Weighted norm comparison, finite clusters, and Fourier escape for exponential polynomials},pdfkeywords={exponential polynomial, confluent frequency, theta weight, Gamma form, coercivity}}
\newtheorem{theorem}{Theorem}[section]
\newtheorem{lemma}[theorem]{Lemma}
\newtheorem{proposition}[theorem]{Proposition}
\newtheorem{corollary}[theorem]{Corollary}
\theoremstyle{definition}
\newtheorem{definition}[theorem]{Definition}
\theoremstyle{remark}
\newtheorem{remark}[theorem]{Remark}
\newcommand{\R}{\mathbb R}
\newcommand{\C}{\mathbb C}
\newcommand{\N}{\mathbb N}
\newcommand{\HH}{\mathcal H_W}
\newcommand{\XX}{\mathcal X_W}
\newcommand{\EE}{\mathcal E_{N,\beta}}
\newcommand{\tg}{\mathfrak t_\Gamma}
\newcommand{\qD}{\mathfrak q_D}
\newcommand{\qth}{\mathfrak q_\theta}
\newcommand{\ip}[2]{\langle #1,#2\rangle}
\newcommand{\norm}[1]{\lVert #1\rVert}
\newcommand{\dd}{\,dx}
\DeclareMathOperator{\Rea}{Re}
\DeclareMathOperator{\Ima}{Im}
\DeclareMathOperator{\supp}{supp}
\allowdisplaybreaks[2]
\setlength{\emergencystretch}{2em}
\title{Uniform Logarithmic Gamma Coercivity\\at Fixed Exponential Complexity}
\author{Jongmin Choi\\[3pt]
\small Independent Researcher, Seoul, Korea\\
\small ORCID: \href{https://orcid.org/0009-0008-7448-514X}{0009-0008-7448-514X}}
\date{1 November 2026}
\begin{document}
\maketitle
\begin{abstract}
Let $W$ be a positive $C^1$ weight such that $W$ and $W'$
decrease faster than every exponential, and let
$P$ range over exponential polynomials of fixed total confluent complexity
whose frequencies lie in a fixed horizontal strip. We prove a uniform
comparison between the weighted norm of $WP$ in $L^2(\R,dx/W)$ and its
ordinary $L^2$ norm, with no separation assumption on the frequencies.
When every real frequency tends in absolute value to infinity, the
normalized Fourier mass of $WP$ escapes every fixed compact interval,
uniformly over the same class. A compact family of constant-coefficient
ordinary differential equations controls collisions; a finite-cluster
argument handles arbitrary real separations. Consequently the exact
logarithmic Gamma form, minus any bounded self-adjoint perturbation in the
weighted space, has Rayleigh quotient tending uniformly to $+\infty$.
We verify the hypotheses for the classical theta weight and prove directly
that its complete prime-power correction is norm-absolutely summable.
An entire two-frequency formula and confluent cosine consequences are
included. Auxiliary examples show why complexity cannot simply be omitted;
they are not actual zeta-zero jets and give no negative-form witness.
No complexity-independent threshold or assertion about all zeta-zero jets
is proved.
\end{abstract}
\noindent\textbf{Keywords.} Exponential polynomials; confluent frequencies; weighted norms;
Fourier escape; logarithmic Gamma form; theta weight.

\section{Introduction and main theorem}

High-frequency modulation makes a logarithmically growing Fourier
multiplier large on any one fixed envelope. Uniformity is subtler when the
envelope is allowed to vary with several complex frequencies, colliding
frequencies, and polynomial factors. The coefficients of a nearly
colliding exponential basis can diverge even when the represented function
remains small. Conversely, a collection of frequencies can split into
widely separated groups, so no one bounded frequency window contains the
whole configuration after a common modulation.

We address these two issues at fixed total complexity. The local estimate
uses the complete annihilating differential equation rather than the
coefficients in an exponential basis. The global argument extracts finitely
many frequency clusters and controls their cross inner products before
normalizing individual envelopes. The resulting constants are uniform in
the real-frequency configuration, but the proof is nonquantitative and
keeps the complexity fixed throughout.

All spaces are complex, and inner products are linear in the first
variable. Ordinary Lebesgue $L^2(\R)$ norms are denoted by $\norm{\cdot}_2$.
We use the unitary Fourier transform
\begin{equation}\label{eq:fourier}
 \widehat u(t)=(2\pi)^{-1/2}\int_\R u(x)e^{-itx}\dd.
\end{equation}

\begin{definition}[Admissible weight]\label{def:weight}
A function $W:\R\to(0,\infty)$ is admissible if it is $C^1$ and, for every
$a>0$, there is $C_a<\infty$ such that
\begin{equation}\label{eq:weight}
 W(x)+|W'(x)|\le C_a e^{-a|x|}\qquad(x\in\R).
\end{equation}
Set $\HH=L^2(\R,dx/W)$, with
$\norm{u}_W^2=\int_\R |u(x)|^2/W(x)\dd$.
\end{definition}

The assumptions imply that $W$ is bounded, has a positive minimum on every
compact interval, and tends to zero at both ends. In particular,
\begin{equation}\label{eq:embedding}
 \norm{u}_2^2\le \norm W_\infty\norm{u}_W^2\qquad(u\in\HH).
\end{equation}
The converse below is specific to a nonlinear class, not to all of $\HH$.

\begin{definition}[Exponential complexity]\label{def:complexity}
For $N\in\N$ and $\beta\ge0$, let $\EE$ consist of zero and all functions
\begin{equation}\label{eq:exp-poly}
 P(x)=\sum_{j=1}^r p_j(x)e^{i\lambda_jx},\qquad
 \sum_{j=1}^r(\deg p_j+1)\le N,\qquad |\Ima\lambda_j|\le\beta,
\end{equation}
where the $\lambda_j\in\C$ are distinct and the $p_j$ are nonzero
polynomials. Repeated frequencies are merged first. Let $\EE(T)$ be the
subclass for which $|\Rea\lambda_j|\ge T$ for every displayed frequency.
\end{definition}

The cost of a degree-$k$ polynomial at one frequency is $k+1$, regardless
of how many of its coefficients vanish. Distinct exponential-polynomial
terms are linearly independent, as recalled in Section~\ref{sec:cosine},
so the merged nonzero representation is unique. These classes need not be
closed under addition at fixed $N$.

For $s>0$ define
\begin{equation}\label{eq:gamma-form}
 \rho(s)=\frac{e^{-s/2}}{1-e^{-2s}},\qquad
 \tg[u]=\int_0^\infty\rho(s)\norm{u(\cdot+s)-u}_2^2\,ds,
 \qquad \XX=\{u\in\HH:\tg[u]<\infty\}.
\end{equation}
The nonnegative integral is interpreted without cancellation. If $D$ is
any bounded self-adjoint operator on $\HH$, set
\begin{equation}\label{eq:q}
 \qD(u,v)=\tg(u,v)-\ip{Du}{v}_W,\qquad
 \qD[u]=\qD(u,u).
\end{equation}
Here $\tg(u,v)$ is the polarized difference form. Every $WP$ in
Definition~\ref{def:complexity} belongs to $\XX$, as proved below.

\begin{theorem}[Uniform finite-complexity coercivity]\label{thm:main}
Fix an admissible weight $W$, $N\in\N$, and $\beta\ge0$.
\begin{enumerate}[label=\textup{(\roman*)},leftmargin=2.1em]
\item There is $C=C(N,\beta,W)<\infty$ such that
\begin{equation}\label{eq:comparison}
 \norm{WP}_W^2\le C\norm{WP}_2^2\qquad(P\in\EE).
\end{equation}
\item For every fixed $0\le R<\infty$,
\begin{equation}\label{eq:escape}
 \varepsilon_{N,\beta,W}(T;R):=
 \sup_{0\ne P\in\EE(T)}
 \frac{\norm{\mathbf1_{[-R,R]}\widehat{WP}}_2}{\norm{WP}_2}
 \longrightarrow0\qquad(T\to\infty).
\end{equation}
\item For every bounded self-adjoint $D$ on $\HH$,
\begin{equation}\label{eq:coercivity}
 \inf_{0\ne P\in\EE(T)}\frac{\qD[WP]}{\norm{WP}_2^2}
 \longrightarrow+\infty,\qquad
 \inf_{0\ne P\in\EE(T)}\frac{\qD[WP]}{\norm{WP}_W^2}
 \longrightarrow+\infty.
\end{equation}
\end{enumerate}
No lower bound on frequency separation or bound on coefficient
conditioning is required. The constants and height thresholds may depend
on $N$, $\beta$, $W$, and, in \textup{(iii)}, $\norm D$.
\end{theorem}

In particular, for every prescribed $L>0$ there is $T_L<\infty$ such that
$\qD[WP]\ge L\norm{WP}_W^2$ for every $P\in\EE(T_L)$.
This is uniform in the represented function norm, not in the Euclidean
norm of raw exponential coefficients. The theorem gives no effective
value of $T_L$ and no uniformity as $N\to\infty$.

Section~\ref{sec:gamma} fixes the exact Gamma form and its domain.
Sections~\ref{sec:ode}--\ref{sec:proof} prove the theorem.
Section~\ref{sec:theta} verifies the theta and full-prime specialization
without importing an energy-realization theorem. The remaining sections
record finite cosine-jet consequences and the limitations of the
complexity hypothesis. No historical-priority claim is made.

\section{The exact Gamma form}\label{sec:gamma}

\begin{lemma}[Multiplier and domain]\label{lem:gamma}
For $u\in\HH$,
\begin{equation}\label{eq:multiplier}
 \tg[u]=\int_\R d_\Gamma(t)|\widehat u(t)|^2\,dt,
 \qquad
 d_\Gamma(t)=2\int_0^\infty\rho(s)(1-\cos(ts))\,ds.
\end{equation}
The multiplier is nonnegative and satisfies
\begin{equation}\label{eq:logbound}
 d_\Gamma(t)\ge\frac12\log(1+4t^2).
\end{equation}
Moreover, for a constant $C_\Gamma<\infty$,
\begin{equation}\label{eq:h1bound}
 \tg[u]\le C_\Gamma\bigl(\norm u_2^2+\norm{u'}_2^2\bigr)
 \qquad(u\in H^1(\R)).
\end{equation}
Thus $WP\in\HH\cap H^1(\R)\subset\XX$ for every exponential polynomial
$P$.
\end{lemma}
\begin{proof}
By~\eqref{eq:embedding}, $u\in L^2$. Plancherel for each shift and Tonelli
for the nonnegative integrand give~\eqref{eq:multiplier}, including when
either side is infinite. Near zero, $\rho(s)=O(s^{-1})$; at infinity,
$\rho(s)=O(e^{-s/2})$. Thus the integral defining $d_\Gamma(t)$ converges
for each $t$. The estimates
\[
 \norm{u(\cdot+s)-u}_2\le s\norm{u'}_2,
 \qquad \norm{u(\cdot+s)-u}_2\le2\norm u_2
\]
prove~\eqref{eq:h1bound} by splitting at $s=1$.

Since $1-e^{-2s}\le2s$,
\[
 d_\Gamma(t)\ge I(t):=\int_0^\infty
 \frac{e^{-s/2}(1-\cos(ts))}{s}\,ds.
\]
Differentiation under the integral on compact $t$ intervals gives
$I'(t)=\int_0^\infty e^{-s/2}\sin(ts)\,ds=t/(t^2+1/4)$.
Integrating from $0$ yields $I(t)=\tfrac12\log(1+4t^2)$.
Finally $P$ and $P'$ grow at most like a polynomial times an exponential.
Equation~\eqref{eq:weight} makes
$\int W|P|^2$, $\int W^2|P|^2$, and
$\int |W'P+WP'|^2$ finite. This proves admission.
\end{proof}

For reference, the exact digamma normalization is
\begin{equation}\label{eq:digamma}
 d_\Gamma(t)=\Rea\psi\left(\frac14+\frac{it}{2}\right)-\psi(1/4),
 \qquad \psi=\Gamma'/\Gamma.
\end{equation}
Indeed, subtract the integral representations
\[
 \psi(z)+\gamma=\int_0^\infty
 \frac{e^{-r}-e^{-zr}}{1-e^{-r}}\,dr\qquad(\Rea z>0)
\]
at $z=1/4+it/2$ and $z=1/4$, take real parts, and set $r=2s$
\cite[Eq.~5.9.16]{DLMF}. This identity names the form; only the directly
proved bound~\eqref{eq:logbound} is needed for coercivity.

\begin{proposition}[Closed form]\label{prop:closed}
The space $\XX$, with norm
$\norm u_X^2=\norm u_W^2+\tg[u]$, is complete. The form $\tg$ is densely
defined and nonnegative on $\HH$. The form $\qD$ is closed and bounded
below by $-\norm D\norm u_W^2$.
\end{proposition}
\begin{proof}
If $u_n$ is Cauchy in $\norm{\cdot}_X$, it converges to $u$ in $\HH$,
hence in ordinary $L^2$. In Fourier variables,
$d_\Gamma^{1/2}\widehat u_n$ is Cauchy in $L^2$. A subsequence of the two
convergent $L^2$ sequences converges almost everywhere, which identifies
the second limit as $d_\Gamma^{1/2}\widehat u$. Therefore $u\in\XX$ and
$u_n\to u$ in $\norm{\cdot}_X$. Compactly supported smooth functions lie
in $\XX$ by Lemma~\ref{lem:gamma} and are dense in $\HH$: truncate first,
then use the positive continuous weight bounds on a compact interval and
ordinary smooth approximation. Finally
$|\ip{Du}{u}_W|\le\norm D\norm u_W^2$ and
\[
 \norm u_X^2\le \qD[u]+(\norm D+1)\norm u_W^2
 \le \tg[u]+(2\norm D+1)\norm u_W^2.
\]
The shifted form norm is equivalent to the complete $X$ norm.
\end{proof}

No ordinary Fourier projection is asserted to preserve $\HH$. Fourier
cutoffs below are used only to estimate integrals in ordinary $L^2$.

\section{Uniform envelopes in a bounded frequency window}\label{sec:ode}

\begin{lemma}[Compact ODE estimate]\label{lem:ode}
Fix $N\in\N$ and $0\le M<\infty$. If $P$ has complexity at most $N$ and all
its frequencies obey $|\lambda_j|\le M$, then, with $A=WP$,
\begin{equation}\label{eq:envelope}
 \norm A_W+\norm{A'}_2\le C_{N,M,W}\norm A_2.
\end{equation}
The estimate is uniform when distinct frequencies approach one another
and when polynomial coefficients vary arbitrarily.
\end{lemma}
\begin{proof}
Fix the total order $m=\sum_j(\deg p_j+1)\le N$. The function $P$
satisfies the monic differential equation
\begin{equation}\label{eq:annihilator}
 \prod_j(\partial_x-i\lambda_j)^{\deg p_j+1}P=0.
\end{equation}
The coefficients of the monic polynomial on the left lie in a compact set:
they are continuous symmetric polynomials in $m$ roots from a closed disk.
Repeated roots are included in that set. Let $B_a$ be the corresponding
companion matrix and set $Y=(P,P',\ldots,P^{(m-1)})^{\mathsf T}$. Then
$Y'=B_aY$, and $\norm{B_a}\le L$ uniformly in this compact family.

We claim the uniform observability estimate
\begin{equation}\label{eq:observable}
 |Y(0)|^2\le C_m\int_{-1}^1|P(x)|^2\dd.
\end{equation}
Otherwise there would be coefficient vectors $a_n$ in the compact family
and unit initial vectors $v_n$ for which the integrals tend to zero.
Extract $a_n\to a$ and $v_n\to v$, with $|v|=1$. The power series of the
matrix exponential converges uniformly for $|x|\le1$ and bounded matrices,
so the first components of $e^{xB_{a_n}}v_n$ converge uniformly to that of
$e^{xB_a}v$. The latter component vanishes identically on $[-1,1]$.
Its derivatives of orders $0,\ldots,m-1$ at zero are the components of $v$,
so $v=0$, a contradiction.

The matrix exponential estimate gives
$|Y(x)|\le e^{L|x|}|Y(0)|$ for all real $x$. The same type of bound holds
for $P'$, including when $m=1$ by its differential equation. Since
$w_*:=\min_{[-1,1]}W>0$,
\[
 \int_{-1}^1|P|^2\dd\le w_*^{-2}\norm{WP}_2^2.
\]
Combining this with~\eqref{eq:observable} yields
\begin{equation}\label{eq:growth}
 |P(x)|+|P'(x)|\le C e^{L|x|}\norm{WP}_2.
\end{equation}
The integrals of $e^{2L|x|}W$, $e^{2L|x|}W^2$, and
$e^{2L|x|}|W'|^2$ are finite by~\eqref{eq:weight}.
Integrating~\eqref{eq:growth} and using $(WP)'=W'P+WP'$ proves the lemma
for order $m$. Take the maximum of the finitely many constants for
$1\le m\le N$.
\end{proof}

The compact family is formed only after frequencies are placed in a
bounded window. The companion matrices of an original configuration with
unbounded real frequencies need not be bounded. No inverse Vandermonde
matrix or lower bound on root spacing is used.

\section{Finite clusters at arbitrary real separations}\label{sec:clusters}

\begin{lemma}[Subsequential cluster decomposition]\label{lem:clusters}
Let $P_n\in\EE$ be nonzero. There is a subsequence, still indexed by $n$,
and a fixed finite collection of clusters $\mathcal C$ such that
\begin{equation}\label{eq:cluster}
 WP_n=\sum_{C\in\mathcal C}e^{it_{C,n}x}A_{C,n},\qquad
 A_{C,n}=WQ_{C,n},
\end{equation}
where each $t_{C,n}$ is the real part of an original frequency,
all frequencies of each $Q_{C,n}$ lie in a fixed bounded disk depending
on $C$, and $|t_{C,n}-t_{D,n}|\to\infty$ for $C\ne D$.
There are constants $L_C<\infty$, independent of $n$, such that
\begin{equation}\label{eq:cluster-envelope}
 \norm{A_{C,n}}_W+\norm{A_{C,n}'}_2\le L_C\norm{A_{C,n}}_2.
\end{equation}
Moreover,
\begin{equation}\label{eq:decouple}
 \norm{WP_n}_2^2=(1+o(1))
 \sum_{C\in\mathcal C}\norm{A_{C,n}}_2^2,
\end{equation}
where the relative error is independent of the cluster amplitudes.
\end{lemma}
\begin{proof}
Extract a subsequence on which the number of distinct frequencies and all
polynomial degrees are fixed; there are only finitely many such choices.
Label the frequencies along this subsequence. For every pair $j,k$, extract
further so that $\Rea\lambda_{j,n}-\Rea\lambda_{k,n}$ converges in the
extended real line. There are finitely many pairs. Declare $j\sim k$ when
this limit is finite. Reflexivity and symmetry are immediate; transitivity
follows by adding the two convergent real differences. The equivalence
classes are the clusters.

Choose a member $j_C$ in each cluster and put
\[
 t_{C,n}=\Rea\lambda_{j_C,n},\qquad
 Q_{C,n}(x)=\sum_{j\in C}p_{j,n}(x)
 e^{i(\lambda_{j,n}-t_{C,n})x}.
\]
Within a cluster the real offsets are bounded, while the imaginary parts
remain in $[-\beta,\beta]$. Between distinct clusters the carrier differences
tend in modulus to infinity. This proves the stated decomposition and,
by Lemma~\ref{lem:ode},~\eqref{eq:cluster-envelope}.

For distinct clusters, integration by parts gives
\begin{align}\label{eq:cross}
 \left|\ip{e^{it_{C,n}x}A_{C,n}}{e^{it_{D,n}x}A_{D,n}}_2\right|
 &\le\frac{\norm{A_{C,n}'}_2\norm{A_{D,n}}_2+
 \norm{A_{C,n}}_2\norm{A_{D,n}'}_2}{|t_{C,n}-t_{D,n}|}\notag\\
 &\le\frac{L_C+L_D}{|t_{C,n}-t_{D,n}|}
 \norm{A_{C,n}}_2\norm{A_{D,n}}_2.
\end{align}
The product $A_{C,n}\overline{A_{D,n}}$ belongs to $W^{1,1}(\R)$ by
Cauchy--Schwarz. Its endpoint limits are zero: the derivative is integrable,
so the limits exist, and its integrability forces those limits to vanish.
Thus the integration by parts is valid.

Write $s=|\mathcal C|\le N$ and $h_{C,n}=\norm{A_{C,n}}_2$. If $s>1$,
let $\delta_n$ be the largest coefficient on the right of~\eqref{eq:cross}.
Then $\delta_n\to0$ and
\[
 \left|\norm{WP_n}_2^2-\sum_C h_{C,n}^2\right|
 \le2\delta_n\sum_{C<D}h_{C,n}h_{D,n}
 \le(s-1)\delta_n\sum_C h_{C,n}^2.
\]
For $s=1$ the identity is exact. This proves~\eqref{eq:decouple}; in
particular, for all sufficiently large $n$,
\begin{equation}\label{eq:amplitude}
 \sum_C\norm{A_{C,n}}_2^2\le2\norm{WP_n}_2^2.
\end{equation}
\end{proof}

The constants $L_C$ depend on the bounded offset windows found along the
subsequence. They are not yet universal constants for the original class.
This is sufficient for the contradiction arguments that follow. Crucially,
\eqref{eq:amplitude} is established without any prior bound on the separate
envelope norms or on the raw coefficients.

\section{Proof of uniform comparison and coercivity}\label{sec:proof}

\begin{proof}[Proof of Theorem~\ref{thm:main}\textup{(i)}]
Suppose no constant in~\eqref{eq:comparison} existed. Choose nonzero
$P_n\in\EE$ with $\norm{WP_n}_W^2>n\norm{WP_n}_2^2$. Apply
Lemma~\ref{lem:clusters}. Multiplication by a real modulation is an
isometry in both weighted and ordinary norms. If $s=|\mathcal C|$ and
$L=\max_C L_C$, the triangle inequality and Cauchy--Schwarz give
\[
 \norm{WP_n}_W^2\le s\sum_C\norm{A_{C,n}}_W^2
 \le sL^2\sum_C\norm{A_{C,n}}_2^2
 \le2sL^2\norm{WP_n}_2^2
\]
for all sufficiently large $n$, a contradiction. Thus a single finite
constant exists for the entire class.
\end{proof}

\begin{proof}[Proof of Theorem~\ref{thm:main}\textup{(ii)}]
If~\eqref{eq:escape} failed, there would be $T_n\to\infty$ and
$P_n\in\EE(T_n)$ such that $\norm{WP_n}_2=1$ and the low-band Fourier
norm remained bounded below by a positive constant. Extract the clusters
of Lemma~\ref{lem:clusters}. Each $t_{C,n}$ is the real part of an original
frequency, so $|t_{C,n}|\ge T_n\to\infty$.

For $|t_{C,n}|>R$, Fourier translation and Plancherel imply
\begin{align}\label{eq:low-band}
 \norm{\mathbf1_{[-R,R]}
 \widehat{e^{it_{C,n}x}A_{C,n}}}_2^2
 &=\int_{|t|\le R}|\widehat A_{C,n}(t-t_{C,n})|^2\,dt\notag\\
 &\le\frac{\norm{A_{C,n}'}_2^2}{(|t_{C,n}|-R)^2}
 \le\frac{L_C^2\norm{A_{C,n}}_2^2}{(|t_{C,n}|-R)^2}.
\end{align}
Equation~\eqref{eq:amplitude} bounds every envelope norm by $\sqrt2$.
There are finitely many clusters, so the triangle inequality makes the
low-band norm of their sum tend to zero, a contradiction.
\end{proof}

\begin{proof}[Proof of Theorem~\ref{thm:main}\textup{(iii)}]
Let $C$ be the constant in~\eqref{eq:comparison} and put
$\varepsilon(T;R)=\varepsilon_{N,\beta,W}(T;R)$. For every fixed $R$,
Lemma~\ref{lem:gamma} and the bounded perturbation estimate give
\begin{equation}\label{eq:master-lower}
 \qD[WP]\ge\left\{
 \frac12\log(1+4R^2)\bigl(1-\varepsilon(T;R)^2\bigr)
 -\norm D C\right\}\norm{WP}_2^2
 \qquad(P\in\EE(T)).
\end{equation}
All Gamma frequencies remain in the original form; its nonnegative
low-frequency part was only discarded in this lower estimate. Given
$L>0$, first choose $R$ so that
$\tfrac14\log(1+4R^2)-\norm D C\ge L$. Then choose $T$ large enough
that $\varepsilon(T;R)^2\le1/2$. This proves the first divergence in
\eqref{eq:coercivity}. For a prescribed weighted coefficient $L_W>0$,
apply that result with $L=L_W C$ and use~\eqref{eq:comparison} to get
$\qD[WP]\ge L_W\norm{WP}_W^2$.
\end{proof}

The order of choices is important: first fix a Fourier band $R$, then
send the minimum real carrier height $T$ to infinity. The extraction proof
does not yield a logarithmic-in-$T$ rate for arbitrary configurations.

\section{The theta weight and complete prime correction}\label{sec:theta}

We now give a concrete realization, including all prime powers. It is
constructed directly on $\HH$; no identification with another energy
space, zero-sum formula, or statement about zeta zeros is used.

For $x\ge0$ define
\begin{equation}\label{eq:theta-kernel}
 \kappa(x)=e^{x/2}\sum_{j=1}^\infty
 (4\tau_j(x)^2-6\tau_j(x))e^{-\tau_j(x)},\qquad
 \tau_j(x)=\pi j^2e^{2x},
\end{equation}
extend evenly, and set
\begin{equation}\label{eq:theta-weight}
 W(x)=\frac{\kappa(x)}{2\cosh(x/2)}.
\end{equation}

\begin{lemma}[Theta estimates]\label{lem:theta}
The function $W$ in~\eqref{eq:theta-weight} is positive, even, smooth, and
admissible. For
\[
 C_\kappa=4\pi^2\sum_{j=1}^\infty j^4e^{-\pi(j^2-1)}<\infty
\]
one has
\begin{equation}\label{eq:theta-tail}
 0<W(x)\le C_\kappa\exp(4|x|-\pi e^{2|x|}).
\end{equation}
Every fixed derivative of $W$ decreases faster than every exponential.
\end{lemma}
\begin{proof}
The series in~\eqref{eq:theta-kernel} and its derivatives converge
uniformly on compact subsets of $[0,\infty)$. Every summand is positive,
since $\tau_j\ge\pi>3/2$. To justify smooth evenness at zero, define on
all of $\R$
\[
 F(x)=e^{x/2}\sum_{j=1}^\infty e^{-\pi j^2e^{2x}}.
\]
This is smooth by locally uniform differentiated convergence. The Gaussian
theta transformation
\begin{equation}\label{eq:jacobi}
 \Theta(t):=\sum_{j\in\mathbb Z}e^{-\pi j^2t}
 =t^{-1/2}\Theta(1/t)\qquad(t>0)
\end{equation}
gives $F(-x)=F(x)+\sinh(x/2)$. For completeness,
\eqref{eq:jacobi} follows by periodizing $y\mapsto e^{-\pi t y^2}$:
its $k$th Fourier coefficient is
$t^{-1/2}e^{-\pi k^2/t}$, by the Gaussian Fourier integral. Both the
periodization and this Fourier series converge absolutely and uniformly;
evaluating at zero proves the identity. This is the usual Jacobi
transformation \cite[Eq.~20.7.32]{DLMF}.

Apply $\partial_x^2-1/4$ to the identity for $F$. Since it annihilates
$\sinh(x/2)$ and commutes with reflection, $F''-F/4$ is even. Direct
differentiation gives precisely~\eqref{eq:theta-kernel} on $x\ge0$.
Thus the even extension is smooth.

For $x\ge0$, dropping the negative term yields
\[
 \kappa(x)\le4\pi^2e^{9x/2}\sum_{j\ge1}j^4e^{-\pi j^2e^{2x}}
 \le C_\kappa e^{9x/2-\pi e^{2x}}.
\]
Since $2\cosh(x/2)\ge e^{x/2}$, this proves~\eqref{eq:theta-tail}.
After $r$ derivatives, each summand of $\kappa$ is $e^{x/2}$ times a
polynomial of degree at most $r+2$ in $\tau_j$, times $e^{-\tau_j}$.
Consequently, for a finite constant $C_r$ and $x\ge0$,
\[
 |\kappa^{(r)}(x)|\le C_r e^{(9/2+2r)x-\pi e^{2x}}.
\]
Here one bounds the remaining sum by a constant multiple of
$\sum_{j\ge1}j^{2r+4}e^{-\pi(j^2-1)}$. Derivatives of
$(2\cosh(x/2))^{-1}$ are bounded by constants times $e^{-x/2}$ on this
half-line. Leibniz's rule, then evenness and smoothness on a compact
interval, proves the derivative assertion and admissibility.
\end{proof}

Write $\Lambda(n)=\log p$ if $n=p^k$ with $p$ prime and $k\ge1$, and
$\Lambda(n)=0$ otherwise. Put
\begin{equation}\label{eq:prime-parameters}
 a_n=\frac{\Lambda(n)}{\sqrt n},\qquad s_n=\log n,
 \qquad c_\Gamma=\log\pi-\psi(1/4).
\end{equation}
Only the elementary bound $0\le\Lambda(n)\le\log n$ will be used.

\begin{proposition}[Norm absolute convergence of all prime powers]\label{prop:primes}
For $s\ge0$, the operator
\[
 (K_su)(x)=W(x)u(x+s)
\]
is bounded on $\HH$, with
\begin{equation}\label{eq:shift-norm}
 \norm{K_s}\le b(s):=\sup_x\sqrt{W(x)W(x+s)}.
\end{equation}
The series
\begin{equation}\label{eq:Db}
 D_\theta=c_\Gamma M_W+
 \sum_{n=2}^\infty a_n(K_{s_n}+K_{s_n}^*)
\end{equation}
converges absolutely in operator norm and defines a bounded self-adjoint
operator on $\HH$. Here $M_Wu=Wu$. Its form is exactly
\begin{align}\label{eq:prime-form}
 \ip{D_\theta u}{v}_W
 &=c_\Gamma\ip uv_2+\sum_{n=2}^\infty a_n
 \left\{\int_\R u(x+s_n)\overline{v(x)}\dd
 +\int_\R u(x)\overline{v(x+s_n)}\dd\right\}.
\end{align}
The complete series in~\eqref{eq:prime-form} is absolutely convergent for
all $u,v\in\HH$.
\end{proposition}
\begin{proof}
With $y=x+s$,
\[
 \norm{K_su}_W^2
 =\int_\R W(y-s)|u(y)|^2\,dy
 \le b(s)^2\norm u_W^2,
\]
which proves~\eqref{eq:shift-norm}. The adjoint is
$(K_s^*u)(x)=W(x)u(x-s)$, as a change of variables in the weighted inner
product shows. Equivalently, weighted Cauchy--Schwarz gives
\begin{equation}\label{eq:corr-bound}
 \left|\int_\R u(x+s)\overline{v(x)}\dd\right|
 \le b(s)\norm u_W\norm v_W.
\end{equation}

At least one of $|x|$ and $|x+s|$ is at least $s/2$. The function
$r\mapsto e^{4r-\pi e^{2r}}$ is decreasing for $r\ge0$, since its
logarithmic derivative is $4-2\pi e^{2r}<0$. Using~\eqref{eq:theta-tail}
for that endpoint and $\norm W_\infty$ for the other gives
\begin{equation}\label{eq:bs}
 b(s)\le(\norm W_\infty C_\kappa)^{1/2}
 \exp\left(s-\frac\pi2e^s\right)\qquad(s\ge0).
\end{equation}
Hence
\begin{equation}\label{eq:prime-summability}
 \sum_{n=2}^\infty a_n b(\log n)
 \le(\norm W_\infty C_\kappa)^{1/2}
 \sum_{n=2}^\infty (\log n)\sqrt n\,e^{-\pi n/2}<\infty.
\end{equation}
The last series converges, for example by the ratio test. Also
$\norm{M_W}\le\norm W_\infty$. Thus~\eqref{eq:Db} converges absolutely
in operator norm, with
\begin{equation}\label{eq:D-norm}
 \norm{D_\theta}\le |c_\Gamma|\norm W_\infty+
 2\sum_{n=2}^\infty a_n b(\log n)<\infty.
\end{equation}
Each partial sum is self-adjoint, so the limit is self-adjoint.
The adjoint formula and~\eqref{eq:corr-bound} prove~\eqref{eq:prime-form}
and its absolute convergence.
\end{proof}

No enumeration of primes, truncation of the prime series, or ordinary
$L^2$ boundedness of the complete prime operator is needed. In this
weighted setting the prime correction is an actual norm-convergent
operator series. Reflection interchanges $K_s$ and $K_s^*$ and commutes
with $M_W$. Thus $D_\theta$ preserves the even subspace. The Gamma form
also preserves the even/odd decomposition because its multiplier is even.

\begin{corollary}[Complete theta and prime form]\label{cor:theta}
For the weight~\eqref{eq:theta-weight}, define
$\qth=\tg-\ip{D_\theta\,\cdot}{\cdot}_W$. For each fixed $N$ and
$\beta$, both Rayleigh quotients in~\eqref{eq:coercivity}, with
$D=D_\theta$, tend uniformly to $+\infty$. The same statement holds
when the class is restricted to even functions. Every prime power and the
exact Gamma normalization are retained.
\end{corollary}
\begin{proof}
Apply Theorem~\ref{thm:main} using Lemma~\ref{lem:theta} and
Proposition~\ref{prop:primes}. Restricting the class cannot lower the infimum.
\end{proof}

\section{Confluent cosine families and finite matrices}\label{sec:cosine}

Let $W$ be even and admissible. For $w\in\C$ and $k\ge0$, put
\begin{equation}\label{eq:jets}
 J_{w,k}(x)=W(x)\left.\partial_z^k\cos(zx)\right|_{z=w}.
\end{equation}
These are even functions; here the word jet means only a parameter
derivative and does not impose a zero equation. Suppose that
$w_1,\ldots,w_r$ are distinct, $\Rea w_j\ge T>0$,
$|\Ima w_j|\le\beta$, and $m_j\ge1$. Set $d=\sum_jm_j$ and
\begin{equation}\label{eq:cosine-combination}
 u=\sum_{j=1}^r\sum_{k=0}^{m_j-1}b_{j,k}J_{w_j,k}.
\end{equation}
The identity
\begin{equation}\label{eq:jet-exp}
 \partial_z^k\cos(zx)=\frac12\bigl((ix)^ke^{izx}+(-ix)^ke^{-izx}\bigr)
\end{equation}
shows that $u/W\in\mathcal E_{2d,\beta}(T)$.

\begin{corollary}[Uniform finite cosine matrices]\label{cor:cosine}
Fix $N$, $\beta$, and a bounded self-adjoint $D$ on $\HH$. For every
$L>0$ there is $T_N(L)<\infty$ such that every combination
\eqref{eq:cosine-combination} with $d\le N$ and
$\min_j\Rea w_j\ge T_N(L)$ satisfies
\begin{equation}\label{eq:cosine-coercivity}
 \qD[u]\ge L\norm u_W^2.
\end{equation}
For the nonredundant prefix list displayed above, the form matrix is
positive definite at those heights. More precisely it dominates $L$ times
the weighted Gram matrix as a Hermitian form on the coefficient space.
\end{corollary}
\begin{proof}
The inequality is Theorem~\ref{thm:main} at complexity $2N$. To justify
strictness of the coefficient form, recall exponential-polynomial
independence. If $\sum_jp_j(x)e^{i\lambda_jx}=0$ with distinct
$\lambda_j$, apply
$\prod_{h\ne j}(\partial_x-i\lambda_h)^{\deg p_h+1}$.
Every term except the $j$th vanishes. After removing $e^{i\lambda_jx}$,
the remaining operator on $p_j$ is a product of powers of
$\partial_x+i(\lambda_j-\lambda_h)$.
Each factor is invertible on the finite polynomial space: in the monomial
basis its matrix is triangular with a nonzero constant diagonal.
Therefore $p_j=0$ for every $j$.

The frequencies $w_1,\ldots,w_r,-w_1,\ldots,-w_r$ are all distinct
because the first group has positive real parts and the $w_j$ are distinct.
Equation~\eqref{eq:jet-exp} and independence imply that the prefix jets
are linearly independent. Hence a nonzero coefficient vector represents a
nonzero $u$, and~\eqref{eq:cosine-coercivity} gives $\qD[u]>0$.
\end{proof}

For repeated or redundant columns the same inequality holds, but the
coefficient matrix can be singular. For a sparse selection at $w_j$, the
correct complexity cost is one plus the largest derivative order selected,
not the number of selected columns. Near collisions the unnormalized Gram
matrix can be ill-conditioned; no uniform lower eigenvalue in raw
Euclidean coordinates is asserted.

This corollary applies to arbitrary frequencies meeting its hypotheses.
It neither uses nor proves that a given frequency is a zero of a special
function. In particular, no conclusion about all actual zeta-zero jets,
a centered half inequality, or the Riemann hypothesis follows from it.

\section{A fixed width cosine packet}\label{sec:packet}

Although the general theorem has no stated rate, a common bounded offset
window permits a useful explicit form of estimate with unevaluated
constants. Fix $R_0\ge0$ and $\beta\ge0$. Let $P_+$ have complexity at most $N$ and frequencies in
$|\Rea\lambda|\le R_0$, $|\Ima\lambda|\le\beta$. Put $A=WP_+$ and
\begin{equation}\label{eq:packet}
 u_T(x)=\frac12\bigl(e^{iTx}A(x)+e^{-iTx}A(-x)\bigr),\qquad T>0,
\end{equation}
where $W$ is even. By Lemma~\ref{lem:ode}, choose a constant $C\ge1$
depending only on $N,R_0,\beta,W$ such that
$\norm A_W,\norm{A'}_2\le C\norm A_2$.

\begin{proposition}[Carrier estimate]\label{prop:packet}
For $T\ge2C$,
\begin{equation}\label{eq:packet-bound}
 \qD[u_T]\ge\left\{
 \frac12\log(1+T^2)\left(1-\frac{16C^2}{T^2}\right)
 -4\norm D C^2\right\}\norm{u_T}_2^2.
\end{equation}
\end{proposition}
\begin{proof}
Integration by parts gives
\[
 \left|\int_\R e^{2iTx}A(x)\overline{A(-x)}\dd\right|
 \le\frac{\norm{A'}_2\norm A_2}{T}
 \le\frac C T\norm A_2^2.
\]
Expanding~\eqref{eq:packet} therefore yields
$\norm{u_T}_2^2\ge\norm A_2^2/4$ for $T\ge2C$.
The weighted triangle inequality gives
$\norm{u_T}_W^2\le C^2\norm A_2^2\le4C^2\norm{u_T}_2^2$.
Using~\eqref{eq:low-band} with the two carriers $\pm T$ and band
$[-T/2,T/2]$, including their factors $1/2$, gives
\[
 \norm{\mathbf1_{[-T/2,T/2]}\widehat{u_T}}_2
 \le\frac{2\norm{A'}_2}{T}
 \le\frac{4C}{T}\norm{u_T}_2.
\]
Now apply~\eqref{eq:logbound} outside that band and subtract
$\norm D\norm{u_T}_W^2$.
\end{proof}

The right side tends to a positive and then arbitrarily large coefficient
as $T\to\infty$ with the window and complexity fixed. The constants have
not been evaluated, so this is not an effective numerical threshold.
If the displayed lower bound is negative at a smaller carrier, that fact
is not a negative value of the form.

\section{Why complexity cannot be removed}\label{sec:limits}

The quantifiers in Theorem~\ref{thm:main} have the order
\[
 \text{for each fixed }N,\beta,L\text{ there is a height }T(N,\beta,L).
\]
They do not provide a single height for all complexities. Taking the span
of all high-frequency jets allows unbounded complexity and is not covered
by the theorem. The following constructions demonstrate that two key
uniform estimates genuinely fail on the larger arbitrary-frequency class.

\begin{proposition}[Failure of complexity independent norm comparison]\label{prop:norm-counter}
For any admissible $W$, let $P_m(x)=x^{2m}$. Then
\begin{equation}\label{eq:norm-counter}
 \frac{\norm{WP_m}_W^2}{\norm{WP_m}_2^2}\longrightarrow\infty.
\end{equation}
\end{proposition}
\begin{proof}
Define the probability measure
\[
 d\nu_m(x)=\frac{W(x)^2x^{4m}\dd}{\int_\R W(y)^2y^{4m}\,dy}.
\]
For each fixed $R\ge1$,
\[
 \nu_m([-R,R])\le
 \left(\frac R{R+1}\right)^{4m}
 \frac{\int_{-R}^R W(x)^2\dd}{\int_{R+1}^{R+2}W(x)^2\dd}
 \longrightarrow0.
\]
The denominator in the fixed ratio is strictly positive. Since
$W(x)\to0$ as $|x|\to\infty$, given $B>0$ choose $R$ so that
$1/W\ge B$ outside $[-R,R]$. Then
\[
 \frac{\norm{WP_m}_W^2}{\norm{WP_m}_2^2}
 =\int_\R\frac1W\,d\nu_m
 \ge B\bigl(1-\nu_m([-R,R])\bigr).
\]
Take the lower limit and then let $B\to\infty$.
\end{proof}

\begin{proposition}[Two high carrier locations with growing degree]\label{prop:escape-counter}
Suppose $W$ is even and admissible. There exist $T_n\to\infty$ and real
even exponential polynomials $P_n$ whose only frequency locations are
$\pm T_n$, but whose complexities are unbounded, such that
\begin{equation}\label{eq:taylor-limit}
 WP_n\longrightarrow W\quad\text{in }\HH\text{ and in }H^1(\R),
 \quad\text{hence in }\XX.
\end{equation}
Consequently uniform low-frequency escape fails when no complexity bound
is imposed. For every fixed bounded self-adjoint $D$,
\begin{equation}\label{eq:q-limit}
 \frac{\qD[WP_n]}{\norm{WP_n}_2^2}
 \longrightarrow\frac{\qD[W]}{\norm W_2^2}<\infty.
\end{equation}
\end{proposition}
\begin{proof}
For fixed $T>0$, put
\[
 S_{m,T}(x)=\sum_{k=0}^m\frac{(-iTx)^k}{k!},\qquad
 P_{m,T}(x)=\frac12\bigl(e^{iTx}S_{m,T}(x)
 +e^{-iTx}S_{m,T}(-x)\bigr).
\]
These functions are real and even on $\R$, have frequencies $\pm T$, and
have complexity $2(m+1)$. For fixed $T$, the Taylor series and its derivative
converge locally uniformly, so $P_{m,T}\to1$ and $P_{m,T}'\to0$ pointwise.
For all $m$,
\[
 |S_{m,T}(x)|\le e^{T|x|},\qquad
 |S_{m,T}'(x)|\le T e^{T|x|},
\]
and hence $|P_{m,T}|\le e^{T|x|}$ and
$|P_{m,T}'|\le2T e^{T|x|}$. Equation~\eqref{eq:weight} supplies integrable
dominating functions for the weighted norm, ordinary norm, and derivative
of $W(P_{m,T}-1)$. Dominated convergence proves
$WP_{m,T}\to W$ in $\HH\cap H^1$ at fixed $T$.

Choose any $T_n\to\infty$ and then choose $m_n$ so large that the sum of
the two norm distances is below $1/n$. Set $P_n=P_{m_n,T_n}$.
Equation~\eqref{eq:h1bound} proves~\eqref{eq:taylor-limit} in $\XX$.
Because $W\in L^1$, $\widehat W$ is continuous and
$\widehat W(0)=(2\pi)^{-1/2}\int W>0$.
Thus some fixed low Fourier band has positive limiting mass, contradicting
any complexity independent escape statement. Finally $\qD$ is continuous
on $\XX$: Cauchy--Schwarz holds for $\tg$, and $D$ is bounded.
This proves~\eqref{eq:q-limit}. The complexities must be unbounded,
otherwise Theorem~\ref{thm:main}\textup{(ii)} would contradict the
nonzero limiting low-band mass.
\end{proof}

These are auxiliary arbitrary-frequency examples. The first has frequency
zero; the second prescribes arbitrarily high polynomial order at two
chosen frequencies. Neither is asserted to use actual zeta zeros or their
actual multiplicities. Neither proves that $\qth$ is negative at any
vector, and~\eqref{eq:q-limit} does not determine the sign of $\qth[W]$.
They rule out deleting the complexity hypothesis from this general
mechanism, not a possible theorem using additional zero-specific structure.

\appendix
\section{An entire two frequency formula}\label{sec:kernel}

For the theta weight and $D_\theta$ of Section~\ref{sec:theta}, define
\begin{equation}\label{eq:overlap}
 e_z(x)=W(x)\cos(zx),\qquad
 H_s(t)=\int_\R W(x+s/2)W(x-s/2)\cos(tx)\dd\quad(s\ge0).
\end{equation}
The following exact formula records the complete form on cosine vectors.
It is not a zero-sum representation and makes no zero-location assumption.

\begin{proposition}[Entire bilinear kernel]\label{prop:kernel}
The function
$\mathcal K(z,v)=\qth(e_z,e_{\overline v})$ is entire on $\C^2$ and is
symmetric in $z,v$, even in each variable, and real under simultaneous
complex conjugation. It satisfies
\begin{align}\label{eq:kernel}
 \mathcal K(z,v)
 &=\int_0^\infty\rho(s)\Bigl[
 H_0(z-v)+H_0(z+v)\notag\\[-2pt]
 &\hspace{3.5em}-\cos\bigl((z+v)s/2\bigr)H_s(z-v)
 -\cos\bigl((z-v)s/2\bigr)H_s(z+v)\Bigr]\,ds\notag\\
 &\quad-\frac{c_\Gamma}{2}\bigl[H_0(z-v)+H_0(z+v)\bigr]\notag\\
 &\quad-\sum_{n=2}^\infty a_n\Bigl[
 \cos\bigl((z+v)s_n/2\bigr)H_{s_n}(z-v)
 +\cos\bigl((z-v)s_n/2\bigr)H_{s_n}(z+v)\Bigr].
\end{align}
The Gamma bracket is kept combined. The integral and series converge
locally uniformly, and every fixed parameter derivative can be taken under
them. In particular,
\begin{equation}\label{eq:jet-matrix}
 \qth(J_{w,k},J_{v,\ell})
 =\left.\partial_z^k\partial_\eta^\ell\mathcal K(z,\eta)
 \right|_{z=w,\,\eta=\overline v}.
\end{equation}
\end{proposition}
\begin{proof}
For a compact parameter set, $e_z$ and each of its fixed parameter
derivatives have uniformly bounded $\HH$ and $H^1$ norms by the theta
tails. The same estimates, with one additional derivative, justify
complex differentiation by dominated convergence in both norms. Thus
$z\mapsto e_z$ is entire into $\XX$ by~\eqref{eq:h1bound}.

Set $C_s(z,v)=\int_\R e_z(x+s)e_v(x)\dd$. With $y=x+s/2$, the weight
product is even in $y$. The product-to-sum identity, whose odd sine terms
integrate to zero, yields
\begin{equation}\label{eq:Cs}
 C_s(z,v)=\frac12\Bigl[
 \cos\bigl((z+v)s/2\bigr)H_s(z-v)
 +\cos\bigl((z-v)s/2\bigr)H_s(z+v)\Bigr].
\end{equation}
In particular $C_s$ is symmetric in $z,v$. Translation invariance gives
\[
 \int_\R(e_z(x+s)-e_z(x))(e_v(x+s)-e_v(x))\dd
 =2C_0(z,v)-2C_s(z,v).
\]
The two prime orientations also both equal $C_{s_n}(z,v)$.
Substitution into the exact forms proves~\eqref{eq:kernel}.

On compact parameter sets, the Gamma bracket is $O(s^2)$ near zero,
bounded by $s^2\norm{e_z'}_2\norm{e_v'}_2$; for large $s$ it is bounded
by $4\norm{e_z}_2\norm{e_v}_2$. Both bounds are integrable against
$\rho(s)\,ds$. They remain valid for fixed parameter derivatives, with
the corresponding differentiated vectors. For the prime part,
\eqref{eq:corr-bound} and~\eqref{eq:prime-summability} give absolute
locally uniform convergence, also after every fixed parameter derivative.
Each $H_s$ is entire by the same theta domination. This proves all
convergence and analyticity assertions. Evenness, symmetry and the real
conjugation identity follow from the formulas; differentiation gives
\eqref{eq:jet-matrix}. Notice that the sesquilinear kernel is
$\qth(e_z,e_v)=\mathcal K(z,\overline v)$, not $\mathcal K(z,v)$.
\end{proof}

\section*{Scope and status}
The proved conclusion is uniform coercivity for each fixed finite
complexity and strip, including collisions and arbitrary real separations.
The full-prime specialization is justified by an operator-norm summability
proof in the weighted space. No positivity at every height, no
complexity-independent threshold, no coefficient-conditioned eigenvalue
bound, and no positivity theorem for all actual zeta-zero jets is claimed.
The manuscript does not invoke a physical energy bridge or a Guinand
zero-sum identity. No numerical research, prime or zero enumeration,
matrix computation, or numerical sign certificate is used. This is a
research manuscript; it is not a claim of peer-reviewed publication or
formal proof-assistant verification.

\begin{thebibliography}{9}
\bibitem{DLMF}
F.~W.~J. Olver, A.~B. Olde Daalhuis, D.~W. Lozier, B.~I. Schneider,
R.~F. Boisvert, C.~W. Clark, B.~R. Miller, B.~V. Saunders, H.~S. Cohl,
and M.~A. McClain, eds.,
\emph{NIST Digital Library of Mathematical Functions}.
\url{https://dlmf.nist.gov/}.
Equations \href{https://dlmf.nist.gov/5.9.E16}{5.9.16} and
\href{https://dlmf.nist.gov/20.7.E32}{20.7.32}; accessed 1 November 2026.
\end{thebibliography}

\medskip
\noindent\small Copyright \textcopyright\ 2026 Jongmin Choi.
This work is licensed under the Creative Commons Attribution-NonCommercial
4.0 International License (CC BY-NC 4.0):
\url{https://creativecommons.org/licenses/by-nc/4.0/}.
\end{document}
